\section{The ledger calculus}
\label{sec:ledger}

Once the transport $x$ is known, a single tail integral contains nearly every
quantity needed by the fitter.

\subsection{The upper-share ledger}

\begin{definition}[Upper-share ledger]
\label{def:ledger}
For an admissible slice define
\begin{equation}
G(z)=\int_z^\infty e^{x(t)}\rho(t)\,\dd t
=\E[Y\1\{Z>z\}].
\label{eq:ledger}
\end{equation}
\end{definition}

The name is literal: $G(z)$ is the fraction of the unit forward carried by
outcomes above log-odds level $z$.  It satisfies
\begin{equation}
G(-\infty)=1,\qquad G(+\infty)=0,\qquad
G'(z)=-e^{x(z)}\rho(z)<0.
\label{eq:ledgerprops}
\end{equation}
For strike $k$, let $z_k$ be the unique root and $u_k$ its rank:
\begin{equation}
x(z_k)=k,\qquad u_k=\Lambda(z_k)=F_X(k).
\label{eq:strikeroot}
\end{equation}

\begin{proposition}[Ledger pricing]
\label{prop:pricing}
For every admissible slice,
\begin{align}
c(k)&=G(z_k)-e^k(1-u_k),
\label{eq:call}\\
P(k)&=e^ku_k-(1-G(z_k)),
\label{eq:put}\\
c'(k)&=-e^k(1-u_k),
\label{eq:cprime}\\
c''(k)&=-e^k(1-u_k)+e^k\frac{\rho(z_k)}{x'(z_k)},
\label{eq:csecond}\\
f_X(k)&=e^{-k}\{c''(k)-c'(k)\}
=\frac{\rho(z_k)}{x'(z_k)}.
\label{eq:bl_division}
\end{align}
\end{proposition}

\begin{proof}
Because $x$ increases, exercise is the event $Z>z_k$.  Splitting asset and
cash legs gives
\[
c(k)=\E[Y\1\{Z>z_k\}]-e^k\Prob(Z>z_k)
=G(z_k)-e^k(1-u_k),
\]
and the put follows on the complementary event.  Differentiating the call,
the moving-root term is
\[
\{G'(z_k)+e^k\rho(z_k)\}\frac{\dd z_k}{\dd k}=0
\]
because $G'(z_k)=-e^{x(z_k)}\rho(z_k)=-e^k\rho(z_k)$.  This yields
\eqref{eq:cprime}.  Since $\dd z_k/\dd k=1/x'(z_k)$, a second derivative
gives \eqref{eq:csecond}; subtracting \eqref{eq:cprime} gives
\eqref{eq:bl_division}.
\end{proof}

Pricing is thus the difference between the share delivered on exercise and
the strike paid on the same event.  The digital probability is $1-u_k$, and
the density is one division of two arrays already present in the transport.
No price differentiation is required in the implementation.

The same rank change of variables prices any integrable terminal payoff:
\begin{equation}
\E[H(Y)]=\int_{-\infty}^{\infty}H(e^{x(z)})\rho(z)\,\dd z.
\label{eq:generalpayoff}
\end{equation}
The ledger is the specialization in which $H(y)=y\1\{y>e^{x(z_0)}\}$ and
the threshold $z_0$ is allowed to vary.  Calls are unusually cheap because
their payoff splits into this asset tail and an ordinary probability tail.
Digitals retain only the probability tail; power and log contracts become
moments of the transport.  The whole payoff calculus therefore remains on one
fixed rank measure.

The constant-speed slice gives a concrete first calculation.  Choose its
scale so that the six-month ATM volatility is exactly twenty percent.  The
solution is $s=\MacToyScale$ and the closed-form shift is $m=\MacToyMu$.
For the call with $k=0.10$, the strike rank is
$u_k=\MacToyOtmPercentile\%$.  The upper-share entry is
$G(z_k)=\MacToyShareOtm$, while the cash entry is
$e^k(1-u_k)=\MacToyCashOtm$.  Hence
\[
c(0.10)=\MacToyShareOtm-\MacToyCashOtm=\MacToyCallOtm,
\]
corresponding to \MacToyIvOtmPct\% implied volatility.  The symmetric law
produces a smile because its exponential return tails carry more remote-strike
value than the lognormal matched at the money.

The fitted law also gives the log-contract variance directly:
\begin{equation}
w_{\rm vs}=-2\E[\log Y]
=-2\int_{-\infty}^{\infty}x(z)\rho(z)\,\dd z.
\label{eq:varswap}
\end{equation}
This is the direct-law form of the familiar strip replication: the classical
construction synthesizes $-\log Y$ from calls and puts over all strikes
\citep{Neuberger1994,DemeterfiDermanKamalZou1999}, while here the law is
already in hand and the same payoff is integrated against its rank measure.
Under the usual continuous-path, continuous-monitoring identity this is the
fair variance-swap strike; jumps and discrete sampling take their standard
separate corrections.

\subsection{The ATM microscope}

At the money, the smile is determined locally by three quantities from the
law.  Let $z_*$ solve $x(z_*)=0$, and write
\begin{equation}
c_0=G(z_*)-(1-u_*),\qquad
u_*=\Lambda(z_*),\qquad
f_*=\frac{\rho(z_*)}{x'(z_*)}.
\label{eq:atm_inputs}
\end{equation}
The ATM Black variance is
\[
w_0=4\left[\PhiN^{-1}\left(\frac{1+c_0}{2}\right)\right]^2,
\qquad d=\frac{\sqrt{w_0}}2,
\]
and define the digital mismatch
\begin{equation}
\Delta=u_*-\PhiN(d).
\label{eq:digitalgap}
\end{equation}

\begin{proposition}[Exact ATM handles]
\label{prop:atm}
The total-variance smile $\Black(k,w(k))=c(k)$ satisfies
\begin{align}
w(0)&=w_0,
&w'(0)&=\frac{2\sqrt{w_0}}{\phin(d)}\Delta,
\label{eq:atm_w}\\
w''(0)&=\frac{2\sqrt{w_0}}{\phin(d)}f_*-2
+\left(\frac1{2w_0}+\frac18\right)w'(0)^2.
\label{eq:atm_wpp}
\end{align}
In volatility units, $\sigma(k)=\sqrt{w(k)/\tau}$,
\begin{align}
\sigma(0)&=\sqrt{w_0/\tau},
&\sigma'(0)&=\frac{\Delta}{\phin(d)\sqrt\tau},
\label{eq:atm_sigma1}\\
\sigma''(0)&=\frac1{\sqrt\tau}
\left[\frac{f_*}{\phin(d)}-\frac1{\sqrt{w_0}}
+\frac{\sqrt{w_0}}4\left(\frac{\Delta}{\phin(d)}\right)^2\right].
\label{eq:atm_sigma}
\end{align}
\end{proposition}

The full differentiation is given in \cref{app:atmproof}.  The first
derivative contains the main idea.  At $(k,w)=(0,w_0)$,
\[
\partial_k\Black=-(1-\PhiN(d)),
\qquad
\partial_w\Black=\frac{\phin(d)}{2\sqrt{w_0}},
\qquad
c'(0)=-(1-u_*).
\]
Implicit differentiation therefore gives
\begin{equation}
w'(0)=\frac{c'(0)-\partial_k\Black}{\partial_w\Black}
=\frac{2\sqrt{w_0}}{\phin(d)}\{u_*-\PhiN(d)\}.
\label{eq:skewderivation}
\end{equation}
The model digital is $1-u_*$; the matched flat-Black digital is
$1-\PhiN(d)$.  Hence $\Delta$ is Black digital minus model digital.  A
negative equity skew is exactly a model digital richer than the flat-Black
digital.  Level, digital, and density determine level, skew, and curvature.

Three sign checks make the interpretation concrete.  For the matched Black
law, $X\sim\mathcal N(-w_0/2,w_0)$, so
$u_*=\Prob(X\le0)=\PhiN(d)$ and the skew vanishes.  If
$\sigma'(0)<0$, then $u_*<\PhiN(d)$ and the model digital exceeds the Black
digital.  Conversely, an upward skew is exactly a cheaper model digital.
The comparison is with $\PhiN(d)$, not with one half: the martingale shift may
place the forward rank on either side of the median even when the skew is
small.

The curvature formula adds one new local fact: $f_*$.  Holding ATM price and
digital fixed, a larger density at the forward raises $w''(0)$.  This matches
the geometry of a peaked distribution: more mass close to the forward makes
nearby options cheap relative to options a little farther away, and the smile
bends upward more rapidly.  The other terms in \eqref{eq:atm_wpp} are the
Black-coordinate correction and a quadratic skew contribution.  Price level,
first CDF value, and first density value form a hierarchy for the three smile
handles.

The identities also run in reverse.  A target ATM level fixes $c_0$; a target
skew fixes $u_*$ through \eqref{eq:atm_w}; and a target curvature fixes $f_*$
through \eqref{eq:atm_wpp}.  These are local constraints on the law rather
than finite differences of a plotted smile.

\subsection{Envelope cancellation}
\label{sec:jacobian}

Calibration requires derivatives with respect to every parameter.  The call
depends on a parameter both through $G$ and through the moving root $z_k$.
The same cancellation used for \eqref{eq:cprime} removes the latter.

\begin{theorem}[Envelope cancellation]
\label{thm:envelope}
For every coordinate $\theta_j$,
\begin{equation}
\frac{\partial c(k;\theta)}{\partial\theta_j}
=\left.\frac{\partial G(z;\theta)}{\partial\theta_j}
\right|_{z=z_k(\theta)}.
\label{eq:envelope}
\end{equation}
\end{theorem}

\begin{proof}
Differentiate \eqref{eq:call}:
\[
\frac{\partial c}{\partial\theta_j}
=\left.\frac{\partial G}{\partial\theta_j}\right|_{z_k}
+\{G'(z_k)+e^k\rho(z_k)\}\frac{\partial z_k}{\partial\theta_j}.
\]
The bracket vanishes by \eqref{eq:ledgerprops} and $x(z_k)=k$.
\end{proof}

This is an envelope theorem: the exercise boundary is the point where the
payoff integrand vanishes, so its displacement has no first-order value.

Because $g$ is affine in the coefficients, all columns are obtained by two
cumulative passes.  With $\bar x=x-m$ and
$I=\int e^{\bar x}\rho$, let $b_j(u)=\partial g(u)/\partial\theta_j$.  Then
\begin{align}
\frac{\partial\bar x(z)}{\partial\theta_j}
&=\int_0^ze^{g(\Lambda(t))}b_j(\Lambda(t))\,\dd t,
\label{eq:pass1}\\
\frac{\partial m}{\partial\theta_j}
&=-\frac1I\int_{-\infty}^\infty
e^{\bar x(z)}\frac{\partial\bar x(z)}{\partial\theta_j}\rho(z)\,\dd z,
\label{eq:pass2}\\
\frac{\partial G(z)}{\partial\theta_j}
&=\int_z^\infty e^{x(t)}
\left(\frac{\partial m}{\partial\theta_j}
+\frac{\partial\bar x(t)}{\partial\theta_j}\right)\rho(t)\,\dd t.
\label{eq:pass3}
\end{align}

The basis columns are fixed functions,
\[
b_L(u)=1-u,\qquad b_R(u)=u,\qquad b_{a_n}(u)=P_n(1-2u).
\]
The first pass differentiates the transport speed and integrates it.  The
normalization dot product removes the component that would change
$\E[e^X]$.  The reverse pass differentiates the upper-share tail.  Every
stage follows the direction of the corresponding primal build, so the
derivative is obtained by array operations rather than a separate
root-and-price routine for each coefficient.

The variance-swap row is obtained from the same sensitivity of $x$:
\begin{equation}
\frac{\partial w_{\rm vs}}{\partial\theta_j}
=-2\int_{-\infty}^\infty
\left(\frac{\partial m}{\partial\theta_j}
+\frac{\partial\bar x(z)}{\partial\theta_j}\right)\rho(z)\,\dd z.
\label{eq:varswaprow}
\end{equation}

Both analytic and finite-difference Jacobians scale as
$O(Pn_{\rm grid})$ for $P=N+1$ parameters, but the analytic route avoids
$P+1$ complete rebuilds.  In the reference implementation it reaches the
same optimum, with costs agreeing to \MacJacSameOptCost\ relative and
parameters to \MacJacSameOptParams, in
\MacJacSpeedupMin--\MacJacSpeedupMax$\times$ less time.  The worst relative
disagreement with an independent central difference is \MacJacFdMaxRel.

The theorem is exact for the continuous curves.  Numerically, $x$ and $G$ are
interpolated separately, so between nodes the price interpolant is not
algebraically identical to the integral of the quantile interpolant.  The
analytic derivative inherits this fourth-order interpolation error.  The
finite-difference comparison in \cref{fig:jacobian} measures the discrete
discrepancy instead of assuming the continuous cancellation transfers
without loss.

\begin{figure}[htbp]
\centering
\paperfig{0.92\textwidth}{fig_jacobian.pdf}
\caption{The analytic Jacobian on the fitted SPY December 2026 slice.  Left:
the normalized price-sensitivity pattern of every parameter across strikes.
Right: columnwise disagreement with an independent central finite difference.
The comparison tests the entire sensitivity pass by a separate route.}
\label{fig:jacobian}
\end{figure}

It is worth counting the shared work.  One build stores $(x,x',G,G')$ on a
fixed grid.  A price is one root and two ledger entries; a density is one
division; the log-contract variance is one weighted sum; the Jacobian reads
ledger sensitivities at the same roots; and calendar order, in
\cref{sec:calendar}, will compare two such arrays.  These are not five
approximations to the law; they are five queries of the same law.
